\documentclass[twoside]{article}
\usepackage[preprint]{aistats2027}
\usepackage[round,authoryear]{natbib}
\usepackage{amsmath,amssymb,amsthm}
\usepackage{mathtools}
\usepackage{microtype}
\usepackage{url}
\renewcommand{\bibname}{References}
\renewcommand{\bibsection}{\subsubsection*{\bibname}}
\bibliographystyle{apalike}
\newtheorem{theorem}{Theorem}
\newtheorem{proposition}[theorem]{Proposition}
\newtheorem{lemma}[theorem]{Lemma}
\newcommand{\E}{\mathbb E}
\newcommand{\R}{\mathcal R}
\newcommand{\sigm}{\operatorname{sigm}}
\begin{document}
\runningauthor{}
\twocolumn[
\aistatstitle{Supplement: The Exact Temperature of Exponential-Weight Aggregation}
\aistatsauthor{Pahan Dewasurendra}
\aistatsaddress{Johns Hopkins University}
]

\section{Exact Risk Identity}

We first prove the algebraic identity used in the lower bound. For any two real predictions $f_1,f_2$, any $w\in[0,1]$, and any label $y$,
\begin{align}
&\{w f_1+(1-w)f_2-y\}^2\nonumber\\
&\quad=w(f_1-y)^2+(1-w)(f_2-y)^2\nonumber\\
&\qquad-w(1-w)(f_1-f_2)^2.
\label{eq:mixtureidentity}
\end{align}
This follows by expanding both sides.

\begin{proof}[Proof of Proposition 2 in the main paper]
On the two-point construction, $\R_{P_n}(f_2)-\R_{P_n}(f_1)=ac/\sqrt n$ and $(f_1-f_2)^2=h^2$ almost surely. Apply \eqref{eq:mixtureidentity} with $w=W_n$, take population expectation conditional on the training sample, and subtract $\R_{P_n}(f_1)$. This gives
\[
\frac{ac}{\sqrt n}(1-W_n)-h^2W_n(1-W_n).
\]
The cumulative empirical loss difference is $aZ_n$, so the definition of the Gibbs weights gives $W_n=\sigm(aZ_n/T)$.
\end{proof}

\section{Triangular Binomial Limits}

Fix $c>0$. For every sufficiently large $n$, let $\epsilon_{n,1},\ldots,\epsilon_{n,n}$ be i.i.d. random signs with
\[
P(\epsilon_{n,i}=1)=\frac12+\frac{c}{2\sqrt n},
\qquad
P(\epsilon_{n,i}=-1)=\frac12-\frac{c}{2\sqrt n},
\]
and put $Z_n=\sum_i\epsilon_{n,i}$. Let $\phi$ and $\Phi$ be the standard normal density and distribution functions.

\begin{lemma}[Global and local limits]
\label{lem:binomial}
The following claims hold.
\begin{enumerate}
\item $Z_n/\sqrt n$ converges in distribution to $N(c,1)$.
\item For each fixed integer $k$ along sample sizes satisfying $n\equiv k\pmod 2$,
\[
\sqrt n\,P(Z_n=k)\longrightarrow2\phi(c).
\]
\item There is $C_c<\infty$ such that, for all sufficiently large $n$ and every integer $k$,
\[
\sqrt n\,P(Z_n=k)\leq C_c.
\]
\end{enumerate}
\end{lemma}

\begin{proof}
The row mean is $\E\epsilon_{n,i}=c/\sqrt n$ and the row variance is $1-c^2/n$. The bounded triangular-array central limit theorem proves the first claim.

For the second and third claims, write $B_n=(n+Z_n)/2$. Then $B_n$ is binomial with success probability $p_n=1/2+c/(2\sqrt n)$. If $k$ has the correct parity,
\[
P(Z_n=k)=P\left(B_n=\frac{n+k}{2}\right).
\]
For fixed $k$, put $d_{n,k}=(n+k)/2-np_n$ and $v_n=np_n(1-p_n)$. Stirling's formula for the binomial mass gives
\begin{align*}
P\left(B_n=\frac{n+k}{2}\right)
&=\frac{1+o(1)}{\sqrt{2\pi v_n}}
\exp\left\{-\frac{d_{n,k}^2}{2v_n}\right\}\\
&=\frac{2\phi(c)+o(1)}{\sqrt n}.
\end{align*}
For the uniform claim, the maximum binomial point mass is at most $C/\sqrt{np_n(1-p_n)}$. Since $p_n(1-p_n)$ is bounded away from zero for all sufficiently large $n$, the required bound follows. Incorrect parity has probability zero.
\end{proof}

We next prove the two limits in the main text. Define
\[
q(u)=\sigm(u)\{1-\sigm(u)\}.
\]
Notice that $0\leq q(u)\leq e^{-|u|}$.

\begin{proof}[Proof of Lemma 3 in the main paper]
The triangular central limit theorem gives $U_n:=Z_n/\sqrt n\Rightarrow U\sim N(c,1)$. For every fixed $\eta>0$,
\[
\sigm(-a\sqrt n U_n/T)
\]
converges uniformly to $\mathbf 1\{U_n<0\}$ on $|U_n|\geq\eta$. The probability of $|U_n|<\eta$ has limit at most $P(|U|\leq\eta)$, which vanishes with $\eta$. Since the variables are bounded, it follows that
\[
\E[1-W_n]=\E\sigm(-aZ_n/T)\longrightarrow P(U<0)=\Phi(-c).
\]

Now restrict to sample sizes of parity $r$. Put $q_{j,r}=q(a(2j+r)/T)$. By Lemma~\ref{lem:binomial}, for every $j\in\mathbb Z$,
\[
\sqrt n\,P(Z_n=2j+r)\longrightarrow2\phi(c).
\]
Moreover,
\[
\sqrt n\,P(Z_n=2j+r)q_{j,r}
\leq C_c e^{-a|2j+r|/T},
\]
and the right side is summable over $j$. Dominated convergence for series gives
\begin{align*}
\sqrt n\,\E[W_n(1-W_n)]
&=\sum_{j\in\mathbb Z}\sqrt n\,P(Z_n=2j+r)
q_{j,r}\\
&\longrightarrow2\phi(c)
\sum_{j\in\mathbb Z}q_{j,r}.
\end{align*}
This is the claimed limit.
\end{proof}

\section{Positive Limits Below Four}

We include all details needed to choose the constants uniformly over both sample parities.

\begin{lemma}[Shifted logistic Riemann sums]
\label{lem:riemann}
For fixed $T>0$, $r\in\{0,1\}$, and $a>0$, put
\[
S_r(a,T)=\sum_{j\in\mathbb Z}q\left(\frac{a(2j+r)}T\right).
\]
Then
\[
\lim_{a\downarrow0}\frac{2a}{T}S_r(a,T)=1.
\]
The convergence holds uniformly over $r\in\{0,1\}$.
\end{lemma}

\begin{proof}
The function $q$ is continuous, nonnegative, and integrable. It is the derivative of $\sigm$, so its integral over the real line is one. The displayed expression is the Riemann sum for $q$ on a grid of spacing $2a/T$. The two values of $r$ correspond to the unshifted grid and its half-grid translate. Both converge to the same integral. Uniformity is immediate because there are only two shifts.
\end{proof}

\begin{lemma}[Gaussian choice]
\label{lem:mills}
For every $T<4$, there is $c>0$ such that
\[
2c\Phi(-c)-\frac{T}{2}\phi(c)>0.
\]
\end{lemma}

\begin{proof}
Integration by parts gives the standard inverse Mills asymptotic
\[
\Phi(-c)=\frac{\phi(c)}c\{1+O(c^{-2})\}.
\]
Therefore $c\Phi(-c)/\phi(c)\to1$. Choose $c$ so that this ratio exceeds $T/4$.
\end{proof}

\begin{proof}[Proof of the lower half of Theorem 1]
Fix $T<4$ and choose $c$ by Lemma~\ref{lem:mills}. For $h\in(0,1)$, let $a=h(2-h)$ and
\[
L_r(h)=ac\Phi(-c)-2h^2\phi(c)S_r(a,T).
\]
Lemma~\ref{lem:riemann} and $a/h\to2$ imply
\begin{align*}
\frac{L_r(h)}h
&=(2-h)c\Phi(-c)
-2h\phi(c)S_r(a,T)\\
&\longrightarrow2c\Phi(-c)-\frac{T}{2}\phi(c)>0.
\end{align*}
Choose one fixed $h>0$ small enough that $L_0(h)>0$ and $L_1(h)>0$. Proposition 2 and Lemma 3 in the main paper show that, along parity $r$,
\[
\sqrt n\,\E[\text{excess risk}]\longrightarrow L_r(h).
\]
With $\kappa=\frac12\min\{L_0(h),L_1(h)\}$, the expected excess is at least $\kappa/\sqrt n$ for every sufficiently large $n$. The same $c,h,f_1,f_2$ work for every such $n$. Only the two input probabilities change.

For a range bound $b>0$, multiply $Y,f_1,f_2$ by $b$. Every squared loss, risk difference, and excess risk is multiplied by $b^2$, while Gibbs weights at temperature $T$ become the normalized weights at temperature $T/b^2$. Hence the lower bound holds whenever $T<4b^2$, with value $\kappa b^2/\sqrt n$. The upper half is Theorem 1 of \citet{hogsgaard2026aggregation}. This completes the exact phase theorem.
\end{proof}

\section{Prior-Weighted Exact Phase}

\begin{proof}[Proof of Corollary 4 in the main paper]
We first prove the upper bound. Let $N=n+1$ and write
\[
S_j=\sum_{i=1}^N(f_j(X_i)-Y_i)^2,
\qquad
p_j=\frac{\pi_je^{-S_j/T}}{\sum_k\pi_ke^{-S_k/T}}.
\]
After observation $i$ is removed, the prior-weighted Gibbs probabilities are
\[
p_j^{(-i)}
=\frac{p_je^{(f_j(X_i)-Y_i)^2/T}}
{\sum_kp_ke^{(f_k(X_i)-Y_i)^2/T}}.
\]
When $T\geq4b^2$, the squared-loss tilting inequality of \citet{hogsgaard2026aggregation} gives
\[
\sum_{i=1}^N
\left\{\sum_jp_j^{(-i)}f_j(X_i)-Y_i\right\}^2
\leq\sum_jp_jS_j.
\]
The Gibbs variational identity and nonnegativity of relative entropy give, for every $j$,
\begin{align*}
\sum_kp_kS_k
&=-T\log\left(\sum_k\pi_ke^{-S_k/T}\right)-T\mathrm{KL}(p\|\pi)\\
&\leq S_j+T\log(1/\pi_j).
\end{align*}
Take expectations. By exchangeability, every term on the left has expectation $\E_{D_n}\R_P(\widehat f_{T,\pi})$. The expected value of $S_j$ is $N\R_P(f_j)$. Dividing by $N$ and minimizing over $j$ proves the stated oracle inequality.

For the converse, it is enough to work at $b=1$. Fix $\pi_1,\pi_2>0$ and put $\lambda=\log(\pi_1/\pi_2)$. The prior-weighted aggregate places weight
\[
W_{n,\pi}=\sigm(aZ_n/T+\lambda)
\]
on $f_1$. Proposition 2 in the main paper is algebraic, so its risk decomposition remains valid with $W_n$ replaced by $W_{n,\pi}$.

Since $\lambda$ is fixed, the triangular-array central limit argument in the proof of Lemma 3 gives
\[
\E[1-W_{n,\pi}]\longrightarrow\Phi(-c).
\]
Along sample sizes of parity $r$, the same local limit and dominated-convergence argument gives
\[
\sqrt n\,\E[W_{n,\pi}(1-W_{n,\pi})]
\longrightarrow2\phi(c)S_{r,\lambda}(a,T),
\]
where
\[
S_{r,\lambda}(a,T)=\sum_{j\in\mathbb Z}q\left(\frac{a(2j+r)}T+\lambda\right).
\]
For every fixed $\lambda$, this shifted lattice satisfies
\[
\frac{2a}{T}S_{r,\lambda}(a,T)\longrightarrow\int_{-\infty}^{\infty}q(u)\,du=1
\qquad(a\downarrow0)
\]
for both parities. The proof of Theorem 1 therefore applies without change: first choose $c$ with $2c\Phi(-c)>(T/2)\phi(c)$, then choose one fixed $h>0$ small enough to make both parity limits positive. Scaling by $b$ completes the proof.
\end{proof}

\section{Sharpness of the Tilting Inequality}

We prove the necessity part of Proposition 5. Sufficiency at $T\geq4$ is Lemma 2 of \citet{hogsgaard2026aggregation}.

Fix $s\in(0,1)$ and let $U$ have probability $1-\varepsilon$ at $s$ and probability $\varepsilon$ at one. Let $p$ denote this law and define its exponential tilt by
\[
q(du)=\frac{e^{u^2/T}p(du)}{\E_p e^{U^2/T}}.
\]
Put $A=e^{s^2/T}$ and $B=e^{1/T}$. Direct differentiation at $\varepsilon=0$ gives
\begin{align}
\E_qU
&=\frac{(1-\varepsilon)sA+\varepsilon B}
{(1-\varepsilon)A+\varepsilon B}\nonumber\\
&=s+\varepsilon\frac BA(1-s)+O(\varepsilon^2)
\label{eq:tiltedexpansion}
\end{align}
and
\begin{align}
(\E_pU^2)^{1/2}
&=\{(1-\varepsilon)s^2+\varepsilon\}^{1/2}\nonumber\\
&=s+\varepsilon\frac{1-s^2}{2s}+O(\varepsilon^2).
\label{eq:rmsexpansion}
\end{align}
The first-order coefficient in \eqref{eq:tiltedexpansion} exceeds that in \eqref{eq:rmsexpansion} exactly when
\[
e^{(1-s^2)/T}>\frac{1+s}{2s},
\]
or equivalently
\begin{equation}
T<A(s):=\frac{1-s^2}{\log\{(1+s)/(2s)\}}.
\label{eq:as}
\end{equation}
In that case $\E_qU>(\E_pU^2)^{1/2}$ for all sufficiently small positive $\varepsilon$, contradicting the tilting inequality.

Let $s=1-d$ and send $d\downarrow0$. Then
\[
1-s^2=2d+O(d^2),
\qquad
\log\frac{1+s}{2s}=\frac d2+O(d^2),
\]
so $A(s)\to4$. For every $T<4$, choose $s$ close enough to one that $T<A(s)$, then choose $\varepsilon>0$ small enough to obtain a violation. Thus universal validity requires $T\geq4$.

For predictions in $[0,b]$, replace $U$ by $bU$. The tilt exponent becomes $b^2U^2/T$, so the normalized temperature is $T/b^2$. The threshold is therefore $4b^2$.

\bibliography{references}
\end{document}
